% =====================================================================
%  PART III -- THE MEASUREMENT  (the only positive result in the paper)
%  Replaces the stub.  English only.  Every number carries its source.
%  Primary sources:
%    research/v3/species/RUN1.md      (beta = 2.5, L_s = 6, 8; run1)
%    research/v3/species/RUN3.md      (beta = 2.5, L_s = 8, N = 1.2e4; run3)
%    research/v3/species/RUN2.md      (estimator identity, m_1 window)
%    docs/position-current.html       (cost origin, four-model ledger)
%    docs/paper-v2/ch12-numerics.tex  (beta = 2.2 volume family)
%    docs/paper-v7/routes20/main.tex  (contact term M_0 ~ a^-8)
%    docs/yang-mills-frontier.md      (gamma_{w_1} = 0)
% =====================================================================

\clearpage
\part{The measurement}

% Section numbering runs continuously through the parts, so that no two
% sections of this paper carry the same number.  (The v3 source restarted
% the counter here, which produced a second \S3, \S4, \S5 and \S6.)
% Cross-references use \label/\ref throughout and are unaffected.

\section{Measuring the atom fraction \texorpdfstring{$w_1$}{w1} at \texorpdfstring{$\beta=2.5$}{beta=2.5}}
\label{sec:measurement}

This part carries the only positive result of the paper, and it is a measurement, not a theorem:
everything in it is graded \LE{} except where a statement is arithmetic or definitional. We define
the quantity (\S\ref{sec:m-what}), report what the lattice actually cost (\S\ref{sec:m-setup}),
give the numbers (\S\ref{sec:m-results}), show that they are determined by the operator and not only
by the theory (\S\ref{sec:m-opdep}), list the three systematics that limit them (\S\ref{sec:m-syst}),
record the two bugs found while producing them (\S\ref{sec:m-bugs}), and say in one place what the
measurement does and does not establish (\S\ref{sec:m-scope}). The order is deliberate: the negative
half of this part is longer than the positive half, and the reason is in \S\ref{sec:m-scope}.

\subsection{What is measured}
\label{sec:m-what}

The object is the atom fraction of the spectral measure of the $0^{++}$ channel at zero momentum.
Fix the operator, its connected two-point function and its spectral measure,
\begin{equation}
O=\tfrac14\operatorname{tr}F^2,
\qquad
C(t)=\langle O(t)O(0)\rangle_c=\int e^{-mt}\,d\nu(m^2),
\qquad
M_0=C(0)=\operatorname{Var}(O),
\label{eq:m-defs}
\end{equation}
let $m_1$ be the lowest mass in $\nu$ and $f_1^2=\nu(\{m_1^2\})=|\langle G|O|\Omega\rangle|^2$ the
weight of the lowest atom. The measured quantity is
\begin{equation}
w_1 \;=\; \frac{f_1^2}{M_0} \;=\; \frac{\nu(\{m_1^2\})}{\nu(\mathbb{R}_+)}\in[0,1],
\qquad
\gamma_{w_1}=0 .
\label{eq:m-w1}
\end{equation}

\begin{claimbox}{cE}
\claimhead{Claim}{(definition, and the one identity that makes it a number).}
For every gauge-invariant $O$ with $\|O\|<\infty$ one has $0\le w_1\le1$; and $w_1$ is invariant
under operator rescaling $O\mapsto ZO$.
\end{claimbox}

\begin{proofx}
$f_1^2=\nu(\{m_1^2\})\le\nu(\mathbb{R}_+)=M_0$ is the statement that one atom is at most the total
mass of a positive measure, and positivity gives $f_1^2\ge0$. For the second part, both $f_1^2$ and
$M_0$ are quadratic in $O$, so $O\mapsto ZO$ multiplies numerator and denominator by $|Z|^2$ and the
ratio is unchanged; hence the anomalous dimension of $w_1$ vanishes, $\gamma_{w_1}=0$. This is the
whole content: the identity fixes the value of $w_1$ under a change of normalisation, and fixes
nothing about its sign. Positivity of $w_1$ is not a scaling statement and is not obtained here.
\end{proofx}

The scaling identity is what makes the number meaningful, and its failure at finite cutoff is the
reason a scheme has to be named before any number is quoted:
\begin{itemize}
\item the numerator is an infrared quantity, a point mass at the lowest atom;
\item the denominator is dominated by the ultraviolet contact term of $O^2$ at coincident points,
  $M_0^{\mathrm{bare}}\sim a^{-8}$;
\item therefore $w_1^{\mathrm{bare}}=f_1^2/M_0^{\mathrm{bare}}\to0$ as $a\to0$. The ratio is dragged
  to zero by the denominator, and the vanishing carries no information about the spectrum.
\end{itemize}
So $w_1$ is not a number but a family $w_1(L,\beta,t)$, with $t$ the renormalisation scheme, i.e.\
the prescription that removes the contact term from $M_0$. Every number in this part belongs to one
explicit member of that family, and to exactly one:
\begin{equation}
w_1 \;=\; \frac{A}{C(0)},
\qquad
C(t)=\frac{1}{L_t}\sum_{t_0}\langle S(t_0+t)S(t_0)\rangle_c,
\qquad
S(t)=\sum_{\mathbf{x}}P(\mathbf{x},t),
\label{eq:m-est}
\end{equation}
with $C(t)$ fitted by $A\,(e^{-mt}+e^{-m(L_t-t)})$ on $t=1,2,3$ at fixed $L_t$, the mass pinned at
$m_1=1.546$, and $P=\tfrac13\sum_{i<j\le3}\tfrac12\operatorname{Re}\operatorname{tr}U_{ij}$ the
zero-momentum spatial plaquette sum. In words: the scheme $t$ here is \emph{the bare operator with
the contact term left in the denominator}, and the measured value is the ratio of the fitted
amplitude to the unsubtracted $t=0$ correlator of the same operator at the same lattice spacing. It
is the only scheme the present data define, and it is the scheme to which $0.3427$ belongs. In
particular it is a finite-$a$ statement: nothing below is a continuum value of $w_1$, and
\S\ref{sec:m-scope} says so again in the negative list.
\src{Definitions and the estimator: \code{eq:m-defs}--\code{eq:m-est} follow the pinned conventions
of \code{RUN1.md:31--39} and \code{RUN3.md:41--47} verbatim (operator, zero-momentum sum, periodic
time averaging, fit window $t=1,2,3$ at both volumes, $w_1=A/C(0)$, $20$-block delete-one jackknife
errors). The pinning mass $m_1=1.546$ is the lower edge of the allowed window derived in
\code{RUN2.md:101--104}. $\gamma_{w_1}=0$ with the reading ``$f_1^2$ and $M_0$ carry the same
$Z_O^2$'' is \code{docs/yang-mills-frontier.md:339}. The contact term
$M_0^{\mathrm{bare}}\sim a^{-8}$ and the resulting $w_1^{\mathrm{bare}}\to0$ are
\code{docs/paper-v7/routes20/main.tex:168--177}.}

\subsection{Lattice setup and cost: how eleven hours became 265 seconds}
\label{sec:m-setup}

The run is four-dimensional $\mathrm{SU}(2)$ Wilson with $\beta=2.5$ ($\kappa=\beta/2=1.25$), spatial
volumes $L_s=6$ and $L_s=8$, temporal extent $L_t=12$, periodic boundary conditions, a
Kennedy--Pendleton heat bath compiled with \code{numba.njit} and \code{fastmath}, and sequential
single-chain updates. The observable is the zero-momentum spatial plaquette sum \eqref{eq:m-est};
one configuration yields the correlator at every operator level at once.

The estimate this run replaces was eleven hours per core for a single volume. The cause was not the
machine and not the algorithm: the plan thermalised each of $10^{5}$ independent configurations
separately, that is, it paid $10^{5}$ times for $400$ equilibration sweeps each. The fix is one line
of logic --- thermalise once per chain, then produce configurations consecutively along that chain,
and measure all operators on each configuration --- and it changes the cost by two orders of
magnitude. Measured:

\begin{table}[htbp]
\centering\small
\begin{tabular}{@{}llrrrr@{}}
\toprule
volume & lattice (sites / links) & $N$ & sweeps/cfg & ms/sweep & wall clock \\
\midrule
$L_s=6$  & $6^3\times12$ \ (2592 / 10368) & $6000$  & $6$ & $2.223$ & $290$ s \\
$L_s=8$  & $8^3\times12$ \ (6144 / 24576) & $4152$  & $3$ & $5.390$ & $265$ s \\
$L_s=8$  & $8^3\times12$, eight chains $\times\,1500$ & $12000$ & $3$ & --- & $34.0$ s \\
\bottomrule
\end{tabular}
\caption{Cost of the measurement. The first two rows are single-threaded, no parallel; the third is
the same volume with eight independent chains over \code{prange}, $352.6$ configurations per second.
Sources: \code{RUN1.md:16--25} (rows 1--2), \code{RUN3.md:20--37} (row 3).}
\label{tab:m-cost}
\end{table}

The correctness check is the plaquette,
$\langle\tfrac12\operatorname{Re}\operatorname{tr}U_p\rangle=0.65270\pm0.00005$ at $L_s=6$,
$0.65236\pm0.00004$ at $L_s=8$, and $0.65234\pm0.00002$ in the independent eight-chain run, all
three consistent with the values tabulated for $\mathrm{SU}(2)$ in \cite{creutz} ($\approx0.65$ at
this coupling). The sampler is therefore
right; the estimator is a separate matter, and \S\ref{sec:m-syst} is about the estimator.

\begin{qbox}
\textbf{The eleven-hour estimate was an implementation problem, not a machine problem.} Eleven hours
per core against $265$ s measured is a factor $1.5\times10^{2}$, and every digit of that factor is
accounted for by re-thermalisation, not by hardware. The same correction is what makes the
statistics affordable later: the run that \S\ref{sec:m-syst} says is still needed
($N\approx4.8\times10^{4}$) costs about $136$ s at the measured throughput, not eleven hours.
\end{qbox}
\src{Cost origin and the ``$10^{5}$ thermalisations'' diagnosis: \code{docs/position-current.html:406}
(``the $11$ h/core of a1 comes from thermalisation being repeated $10^{5}$ times''). Plaquettes and
sweep times: \code{RUN1.md:23--25}, \code{RUN3.md:20--21,33}. The factor
$1.5\times10^{2}=11\,\mathrm{h}/265\,\mathrm{s}$ is arithmetic on these two numbers.}

\subsection{Results}
\label{sec:m-results}

\begin{table}[htbp]
\centering\small
\begin{tabular}{@{}llrrrr@{}}
\toprule
operator & volume & $w_1$ (at $m_1=1.546$) & $w_1/\sigma$ & $N$ & run \\
\midrule
raw ($\tfrac14\operatorname{tr}F^2$) & $L_s=6$, $L_t=12$ & $0.373\pm0.018$  & $20.7\sigma$ & $6000$  & run1 \\
raw ($\tfrac14\operatorname{tr}F^2$) & $L_s=8$, $L_t=12$ & $0.342\pm0.015$  & $22.8\sigma$ & $4152$  & run1 \\
raw ($\tfrac14\operatorname{tr}F^2$) & $L_s=8$, $L_t=12$ & $0.3427\pm0.0119$ & $28.8\sigma$ & $1.2\times10^{4}$ & run3 \\
\bottomrule
\end{tabular}
\caption{The measurement. Errors are $20$-block delete-one jackknife. Sources:
\code{RUN1.md:47,52} (rows 1--2), \code{RUN3.md:51} (row 3); the $w_1/\sigma$ column is the ratio of
the two columns to its left ($0.373/0.018=20.7$, $0.342/0.015=22.8$, $0.3427/0.0119=28.8$).}
\label{tab:m-results}
\end{table}

Three readings of Table~\ref{tab:m-results}, in increasing order of how much they cost:

\begin{enumerate}
\item \emph{Reproducibility.} The two $L_s=8$ runs are independent: different seeds, $8$ chains
  instead of one, $N$ larger by a factor $2.9$. They agree to $+0.0004$, that is $0.02\sigma$. The
  number did not move when the statistics tripled, and that is the strongest single statement in
  this part.
\item \emph{Significance at fixed mass.} $0.3427/0.0119=28.8\sigma$ away from zero. This is the
  significance of $w_1$ with $m_1$ pinned; it is \emph{not} the signal-to-noise of the correlator,
  which is $5.0$ and is discussed in \S\ref{sec:m-syst}. The two numbers are not interchangeable and
  we do not present the larger one as if it were the smaller one's replacement.
\item \emph{Volume dependence.} $(0.373-0.342)/0.373=8.3\%$ between $L_s=6$ and $L_s=8$: the
  fraction does not dilute with volume, in contrast with the free-field behaviour.
\end{enumerate}

The volume trend is available at a second coupling, where the same estimator was run on three
volumes: at $\beta=2.2$ the unsmeared ratio runs $C_1/C_0=0.1367\to0.1229\to0.1216$ for $L=4,6,8$,
a drop and then a plateau rather than a $1/L$ decay, and the three-volume fit of the upper bound
$\hat w_1(t)=e^{mt}C(t)/C(0)$ to $w_\infty+c/L$ gives $w_\infty\approx0.54$. The worst point of that
family is therefore at $L=4$, the smallest volume, and not at $L=\infty$: the finite volume pushes
the fraction \emph{down}, so the $\beta=2.5$ pair sits inside a plateau rather than on a falling
edge. This is a trend, not a limit, and it is the reason we do not extrapolate
$w_1(2.5)$ to $L=\infty$ anywhere in this paper.
\src{Agreement $0.02\sigma$ and the $28.8\sigma$ ratio: \code{RUN3.md:11--12,51}. The $8.3\%$:
\code{RUN1.md:59} and \code{docs/paper-v2/ch12-numerics.tex:281}. The $\beta=2.2$ family
$0.1367\to0.1229\to0.1216$ for $L=4,6,8$, the fit to $w_\infty+c/L$ with $w_\infty\approx0.54$, and
the reading ``inside the plateau, not on a falling edge'': \code{docs/paper-v2/ch12-numerics.tex:284--288};
the same family and the statement that the infimum sits at $L=4$: \code{docs/position-current.html:382}.}

\subsection{Operator dependence}
\label{sec:m-opdep}

This subsection is the most important one in the paper, because it converts the number above from a
fact about Yang--Mills into a fact about a pair.

\begin{table}[htbp]
\centering\small
\setlength{\tabcolsep}{5pt}
\begin{tabular}{@{}lrrrr@{}}
\toprule
basis (same chain, same $12000$ configurations) & $C(0)$ & $C(1)/C(0)$ & $w_1$ & $S/N(2)$ \\
\midrule
raw ($\tfrac14\operatorname{tr}F^2$)      & $19.888$ & $0.07372\pm0.00265$ & $0.3427\pm0.0119$ & $5.0$ \\
spatial APE, $n=1$                        & $45.527$ & $0.00302\pm0.00297$ & $0.0168\pm0.0142$ & $0.9$ \\
spatial APE, $n=2$                        & $39.428$ & $0.00150\pm0.00298$ & $0.0107\pm0.0136$ & $1.6$ \\
\bottomrule
\end{tabular}
\caption{One Markov chain, one set of $12000$ configurations, three operators, all at $m_1=1.546$.
Only the operator changes. Source: \code{RUN3.md:49--53}.}
\label{tab:m-bases}
\end{table}

The ratio of the first row to the second is $0.343/0.017\approx20.2$. Same chain, same
configurations, same jackknife blocks, same fit window: nothing statistical distinguishes the two
numbers, and the difference is not a bug (\S\ref{sec:m-bugs} lists the bugs that were found, and this
is not one of them).

\begin{v3box}
\textbf{The twenty-fold difference is a theorem about notation, realized on data.} Take a
Hamiltonian $H$ whose lowest atom $m_1^2$ is two-fold degenerate, with orthonormal vectors
$\psi_1,\psi_2$, and fix the source family
\[
O_\theta=\mathbf 1+\varepsilon\,(\cos\theta\,\psi_1+\sin\theta\,\psi_2),
\qquad \varepsilon\ll1 .
\]
$H$ is untouched, so the spectrum is untouched: $m_1$, every $m_n$, and the support of $\nu$ are the
same for every $\theta$. But the overlap of the excited part with the lowest atom is
$\varepsilon\cos\theta$, so $f_1^2=\varepsilon^2\cos^2\theta$, while in the scheme $t$ that removes
the contact term carried by $\mathbf 1$ one has $M_0=\varepsilon^2$. Hence
\[
w_1=\cos^2\theta\in[0,1],
\qquad \text{spectrum fixed.}
\]
As $\theta$ runs over $[0,\pi/2]$, $w_1$ sweeps the whole unit interval at a fixed, massive,
unchanged spectrum. The statement is \LT{} --- the proof is \S\ref{sec:reduction-operator} --- and on
these configurations it is additionally \LE, since the factor $20.2$ of Table~\ref{tab:m-bases} is
its instance at $\cos^2\theta=0.049$.
\end{v3box}

The consequence is a statement about language, and it is not negotiable: the sentence
``$w_1(2.5)>0$'' has no content until $O$ is specified. $w_1$ is a property of the pair $(H,O)$ and
not of $H$ alone; Table~\ref{tab:m-bases} is that sentence instantiated on real configurations, with
the factor $20.2$ standing in for the free parameter $\theta$. Two corollaries follow.
\begin{itemize}
\item A change of operator basis is not a change of statistics, so no amount of $N$ will move
  $0.0168$ towards $0.3427$; and no amount of $N$ will make the two agree.
\item The uniqueness statement \texttt{N09-1} --- ``any quantity that passes the four-model test is a
  function of $m$'' --- is true only under \textbf{fixed $O$}. Under a change of $O$ at fixed $H$ the
  quantity moves by a factor $20$ while $m$ does not move at all. This is ledger item $38$, and it is
  the reason the qualifier is now written into the statement rather than left implicit. \LC{}
\end{itemize}
\src{Table~\ref{tab:m-bases} and the factor $20.2$: \code{RUN3.md:49--53}; the same factor quoted as
``$0.343$ vs $0.017$, factor $20$'' at \code{docs/position-current.html:378}. The
$\cos^2\theta$ witness is the one-line counterexample of route \code{u11}, quoted in the abstract of
this paper; the ledger entry is item $38$ at
\code{docs/paper-v7/ledger/main.tex:268}, with the fixed-$O$ qualifier at
\code{docs/position-current.html:378}.}

\subsection{Three systematics}
\label{sec:m-syst}

None of the three is fatal and all three are load-bearing for \S\ref{sec:m-scope}; we state them in
the order in which they would mislead a reader who stopped after Table~\ref{tab:m-results}.

\begin{warnbox}
\begin{enumerate}[label=\textbf{\arabic*.},leftmargin=2.1em]
\item \textbf{The signal-to-noise is below the campaign's own threshold, and the $1/\sqrt N$ route to
  it is dead.} $S/N(2)=C(2)/\sigma_{\mathrm{JK}}=5.0$, with
  $C(2)/C(0)=0.01189\pm0.00238$. The extrapolation used in the planning stage --- $\sigma\propto
  N^{-1/2}$, so $5.9$ at $N=4152$ should become $10$ at $N=1.2\times10^{4}$ --- failed: $N$ rose by
  $2.9\times$ and $S/N$ went from $5.9$ to $5.0$. The requirement is therefore
  $N\approx1.2\times10^{4}\times(10/5.0)^{2}\approx4.8\times10^{4}$, about $136$ s at the measured
  $352.6$ configurations per second. Statistics can be bought; they were not bought here, and the
  threshold $S/N\ge10$ is not met.
\item \textbf{The time correlation is killed by the operator basis, so the GEVP route is unavailable.}
  With spatial APE ($n=1$, $\alpha=0.5$, time direction not smeared, staple normalised $/4$),
  $C(1)/C(0)=0.0030\pm0.0030$ and $C(2)/C(0)=+0.0030\pm0.0032$; $n=2$ gives
  $0.0015\pm0.0030$ and $+0.0034\pm0.0022$. The earlier $/6$ normalisation of run1 gives the same
  $10^{-3}$. So $C(1)\to0$ with $S/N\le1.6$: the projected correlator is a contact term, not a
  propagating state, and a generalised eigenvalue extraction of $m_1$ --- the standard variational
  spectroscopy of \cite{morningstar} --- hence any statement about
  the $\beta$- or $a$-dependence of $w_1$, cannot be built on this basis at $L\le8$.
\item \textbf{The dominant uncertainty is the smearing level, not the volume.} The spread induced by
  the smearing parameter $s$ is $0.29$, against $0.07$ for the volume $L$: a factor $4$. Since the
  $L$-dependence of Table~\ref{tab:m-results} is $8.3\%$, quoting it as an error bar while holding the
  operator fixed would understate the operator dependence of \S\ref{sec:m-opdep} by the very factor
  the quantity is most sensitive to.
\end{enumerate}
\end{warnbox}
\src{(1) $S/N(2)=5.0$, $C(2)/C(0)=0.01189\pm0.00238$, the failed $1/\sqrt N$ extrapolation
($5.9\to5.0$ at $2.9\times$ the statistics), and $N\approx4.8\times10^{4}$:
\code{RUN3.md:13--14,51,84}; the $136$ s is $4.8\times10^{4}/352.6$ on the throughput of
\code{RUN3.md:37}. Note one internal inconsistency in that file, which we resolve in favour of the
table and flag rather than repair: the prose at \code{RUN3.md:14} writes $\pm0.0038$, which would
give $S/N(2)=3.1$ and contradicts its own quoted $5.0$; the table value $\pm0.00238$ reproduces
$0.01189/0.00238=5.0$ and is the one used here. (2) The three APE levels, the $/4$ versus $/6$
normalisation, $S/N(2)\le1.6$ and ``this GEVP route is unusable'': \code{RUN3.md:8--10,52--53,56,66}.
(3) ``$s$-broadening $0.29$ is $4\times$ the $L$-broadening $0.07$'': \code{U14.md:88--89} as
recorded at \code{RUN2.md:8}.}

\subsection{Two bugs found}
\label{sec:m-bugs}

Both are recorded because the second one is the mechanism by which a wrong number would have entered
the paper, and the first one is the mechanism by which a whole class of earlier numbers may already
have.

\begin{qbox}
\textbf{Bug 1: smearing mutated the Markov chain in place.} The first version of the smeared run
passed \code{Ut = U[c]} --- a \code{numpy} \emph{view}, not a copy --- into the APE routine, which
modified it in place. The chain was then continued from the smeared configuration instead of the
sampled one. Symptoms: heat-bath acceptance degraded as $x\to\pm1$, and a single job went from about
$40$ s to more than $10$ min. Fix: allocate \code{buf = np.empty\_like(U[c])} and copy per
configuration. Consequence: \textbf{any earlier result of the form ``smear, then re-measure'' may be
contaminated}, and we do not certify a smeared level produced before the fix. The raw ($n=0$) numbers
of Tables~\ref{tab:m-results} and~\ref{tab:m-bases} are unaffected, since the raw path applies no
smearing at all.

\textbf{Bug 2: the self-contradiction that was not one.} The run chain reported $w_1=0.577$ against
the strict ceiling $\hat w_1(1)=e^{m_1}C(1)/C(0)=0.551$. Since $w_1\le\hat w_1(1)$ is an identity,
the pair $(0.577,\,0.551)$ is impossible, and it was flagged as an internal contradiction. It is not:
the mass had been pinned at $m=1.50$, and the reported value is consistent only for $m_1\ge1.546$.
At $m_1=1.546$ the same data give $0.373\pm0.018$. The resolution is therefore that \emph{the pinned
value was wrong}, and the identity is exactly what caught it. Two consequences we keep: every $w_1$
in this part is quoted at $m_1=1.546$, the lower edge of the allowed window, so they are
\emph{lower edges and not central values}; and ``report $\hat w_1(t)$ alongside $w_1$'' is now a
mandatory step of the analysis, not a courtesy.
\end{qbox}
\src{Bug 1: \code{RUN3.md:76--78}. Bug 2: the ceiling $\hat w_1(1)=e^{m_1}C_1/C_0=0.551$, the
reported $0.577$, the window $m_1\ge1.546$ and the instructive reading of the pair:
\code{RUN2.md:101--104}; the pinned-$1.546$ values are those of Table~\ref{tab:m-results}.}

\subsection{What this does and does not establish}
\label{sec:m-scope}

\begin{v3box}
\textbf{What can be said.} For the raw operator $O=\tfrac14\operatorname{tr}F^2$, at $\beta=2.5$,
$L_s=8$, $L_t=12$, in the scheme \eqref{eq:m-est}, the atom fraction is
\[
w_1=0.3427\pm0.0119 ,
\]
which is $28.8\sigma$ away from zero and reproduced by an independent run at $0.02\sigma$; the
dependence on the volume between $L_s=6$ and $L_s=8$ is $8.3\%$; and at $\beta=2.2$ the family over
$L=4,6,8$ decreases and then flattens, with the infimum at the smallest volume. That is the positive
content of this paper, and it is a measurement (\LE{}), not a floor.
\end{v3box}

\begin{warnbox}
\textbf{What cannot be said.} Four sentences, each of which the numbers above forbid.
\begin{enumerate}[label=\textbf{(\roman*)},leftmargin=2.3em]
\item \emph{``$w_1>0$ for any operator.''} A change of basis on the same configurations gives
  $0.0168\pm0.0142$, a factor $20.2$ smaller (\S\ref{sec:m-opdep}). $w_1$ is a property of $(H,O)$.
\item \emph{``The signal-to-noise is enough.''} $S/N(2)=5.0<10$ (\S\ref{sec:m-syst}, item 1). The
  $28.8\sigma$ of Table~\ref{tab:m-results} is the significance of $w_1$ at pinned mass and does not
  replace it.
\item \emph{``Time correlations are present.''} $C(1)/C(0)=0.003$ and the GEVP route is unavailable
  on these volumes (\S\ref{sec:m-syst}, item 2).
\item \emph{``This advances Clay.''} $w_1$ and $m$ do not imply one another in either direction:
  \begin{itemize}
  \item \textbf{2D Yang--Mills:} $m=3g^2/8>0$ while $\nu_{\frac14\operatorname{tr}F^2}\equiv0$; and
    more strongly, $m(L)=\sigma L\to\infty$ while $w_1=0$. A local operator is blind to the gap.
  \item \textbf{M6} (oscillator $\otimes$ free module): $w_1=0.5$ while $m=0$.
  \item \textbf{Free field:} $w_1=0$ while $m=0$.
  \end{itemize}
  So this measurement does \textbf{not} prove the mass gap, and the direction of the failure is not
  the direction a reader might expect: the failure is not that the number is small or noisy, it is
  that the number is orthogonal to the gap.
\end{enumerate}
\end{warnbox}

One qualification about the fourth item, because it is easily confused with \S\ref{sec:m-opdep}, and
the two must be kept apart rather than merged. \S\ref{sec:m-opdep} varies $O$ at \emph{fixed} $H$:
one theory, two operators, a factor $20.2$. Item (iv) compares \emph{different} $H$: four models, in
each of which $w_1$ and $m$ point in opposite directions. The first says the quantity is not a
function of the theory alone; the second says that even after the operator is fixed, the quantity is
not a function of the gap. Both are needed, and neither implies the other.

\medskip
\noindent\textbf{What the residual actually is.} Not statistics: \S\ref{sec:m-syst} item 1 says the
statistics cost $136$ s and can be bought. Not the volume: item 3 says the volume contributes $0.07$
against $0.29$ for the operator. What is missing is an operator basis that keeps a propagating
correlation at these volumes (item 2), and --- above that --- an argument that connects the atom
fraction to the mass at all, which item (iv) shows cannot be an implication in either direction. The
measurement is honest, reproducible, and it does not move the verdict.
\src{Item (iv) in full, including $m=3g^2/8>0$ with $\nu_{\frac14\operatorname{tr}F^2}\equiv0$, the
stronger $m(L)=\sigma L\to\infty$ with $w_1=0$, M6 with $w_1=0.5$ and $m=0$, and the free field with
$w_1=0$ and $m=0$: \code{docs/position-current.html:392}. The cost claim is
Table~\ref{tab:m-cost} and the $136$ s arithmetic of \S\ref{sec:m-syst}; the $0.07$ versus $0.29$
split is item 3 of \S\ref{sec:m-syst}. The grading \LE{} for the measurement and the refusal to call
it a floor follow the provenance rule of the front matter.}
